% Part: set-theory
% Chapter: story
% Section: predicativity

\documentclass[../../../include/open-logic-section]{subfiles}

\begin{document}

\olfileid{sth}{story}{predicative}

\olsection{প্রেডিকেটিভ ও ইমপ্রেডিকেটিভ}

রাসেলের সেট~$R$-কে $\Setabs{x}{x \notin x}$ দিয়ে সংজ্ঞায়িত করা
হয়েছিল। আরও স্পষ্ট করে বললে, $R$ সেই সেট, যার সদস্য ঠিক সেইসব
% BN-SRC-430: the source's “not non-self-membered” reverses x not-in x.
সেট, যেগুলি নিজেদের সদস্য নয়। তাই $R$-কে সংজ্ঞায়িত করার সময়
আমরা এমন একটি ক্ষেত্রের উপর পরিমাণায়ন করি, যেখানে $R$ নিজেও
থাকত (যদি সেটি থাকত)।

এটি একটি \emph{ইমপ্রেডিকেটিভ} সংজ্ঞা। আরও সাধারণভাবে, কোনো
সংজ্ঞায় যে ক্ষেত্রের উপর পরিমাণায়ন করা হয়, সংজ্ঞায়িত বস্তুটি
নিজেই সেই ক্ষেত্রে থাকলে, সংজ্ঞাটিকে ইমপ্রেডিকেটিভ বলা যায়।

কূটাভাসগুলি প্রকাশিত হওয়ার পরে হোয়াইটহেড, রাসেল, পুয়াঁকারে
ও ভাইল এই ধরনের ইমপ্রেডিকেটিভ সংজ্ঞাকে ``দুষ্টচক্রপূর্ণ'' বলে
প্রত্যাখ্যান করেছিলেন:
\begin{quote}
  এড়াতে হবে এমন কূটাভাসগুলি বিশ্লেষণ করলে দেখা যায়, সবগুলিই
  এক ধরনের দুষ্টচক্র থেকে আসে। সংশ্লিষ্ট দুষ্টচক্র তৈরি হয়
  এই অনুমান থেকে যে, একটি বস্তুর সংগ্রহে এমন সদস্য থাকতে পারে,
  যাদের কেবল সমগ্র সংগ্রহটির সাহায্যেই সংজ্ঞায়িত করা যায়
  [\ldots. \textparagraph]

  অবৈধ সমগ্রতা এড়ানোর নীতিটি এভাবে বলা যায়: ``কোনো সংগ্রহের
  \emph{সবগুলি} নিয়ে যার সংজ্ঞা, সেটি সেই সংগ্রহের সদস্য হতে পারে
  না''; অথবা বিপরীত দিক থেকে: ``কোনো সংগ্রহকে সমগ্র হিসেবে
  ধরলে তাতে এমন সদস্য থাকত, যাদের শুধু সেই সমগ্রটির সাহায্যে
  সংজ্ঞায়িত করা যায়—এমন হলে ওই সংগ্রহের কোনো সমগ্র নেই।''
  অবৈধ সমগ্রতা ধরে নেওয়ার সঙ্গে জড়িত দুষ্টচক্র এড়াতে এই
  নীতি সাহায্য করে বলে আমরা একে ``দুষ্টচক্র নীতি'' বলব।
  \citep[p.~37]{WhiteheadRussell1910}
\end{quote}
% BN-SRC-431: the source says rejecting the principle, reversing the argument.
তাদের অনুসরণ করে যদি আমরা এই \emph{দুষ্টচক্র নীতি} গ্রহণ করি,
তবে \olref[sth][story][rus]{sec}-এর বিপর্যয়কর অনানুষ্ঠানিক
ধর্মনির্দেশে সেট-গঠন ছকের বদলে এমন কিছু গ্রহণের চেষ্টা করতে পারি:

\begin{defish}
\emph{প্রেডিকেটিভ ধর্মনির্দেশে সেট-গঠন।} কেবল সেটের উপর
পরিমাণায়ন করে এমন প্রত্যেক সূত্র~$\phi$-এর জন্য
সেট$^\prime$ $\Setabs{x}{\phi(x)}$-এর অস্তিত্ব আছে।
\end{defish}

যতক্ষণ সেট$^{\prime}$-গুলি সেট না হয়, ততক্ষণ কোনো স্ববিরোধ
উঠবে না।

দুর্ভাগ্যবশত, প্রেডিকেটিভ ধর্মনির্দেশে সেট-গঠন খুব বেশি কিছু
\emph{গঠন} করতে পারে না। এটি নতুন সত্তা—সেট$^\prime$—যোগ করে।
তাই সেট$^\prime$-এর উপর পরিমাণায়ন করে এমন সূত্রও বিবেচনা
করতে হবে। সেগুলি সব সময় সেট$^\prime$-ই দিলে নিজেদের সদস্য
নয় এমন সব সেট$^\prime$-এর সেট$^\prime$ বিবেচনা করলেই
রাসেলের কূটাভাস আবার দেখা দেবে। একই চিন্তাধারা অনুসরণ করে
তাই বলতে হবে, সেট$^\prime$-এর উপর পরিমাণায়নকারী সূত্র
সংশ্লিষ্ট একটি সেট$^{\prime\prime}$ দেয়। তারপর দরকার হবে
সেট$^{\prime\prime\prime}$, সেট$^{\prime\prime\prime\prime}$,
ইত্যাদি। এত প্রাইমচিহ্নের বদলে এগুলিকে সেট$_0$, সেট$_1$,
সেট$_2$, সেট$_3$, সেট$_4$,\ldots{} হিসেবে ভাবা সহজ। এভাবে
আমরা (সরল) প্রকারতত্ত্বের দিকে এগোতে পারি।

এ ধরনের তত্ত্বের বিরুদ্ধে কয়েকটি আপাত আপত্তি আছে (যদিও
সেগুলি \emph{নির্ণায়ক} কি না, তা স্পষ্ট নয়)। সংক্ষেপে:
এই তত্ত্ব ব্যবহার করা জটিল; এটি নানা ধরনের পৃথক বস্তুর অস্তিত্ব
ঢালাওভাবে ধরে নেয়; আর ইমপ্রেডিকেটিভ সংজ্ঞা শেষ বিচারে
\emph{সত্যিই এত খারাপ} কি না, তা-ও স্পষ্ট নয়।

শেষের কথাটি বোঝাতে \citeauthor{Ramsey1925}-র একটি মন্তব্য দেখো:
\begin{quote}
  কোনো দলের সবচেয়ে লম্বা মানুষ বলে আমরা কাউকে নির্দেশ করতে
  পারি; এভাবে তিনি নিজে যে সমগ্রটির সদস্য, সেটির সাহায্যেই
  তাঁকে শনাক্ত করা হয়, অথচ তাতে কোনো দুষ্টচক্র তৈরি হয় না।
  \citep{Ramsey1925}
\end{quote}
র‍্যামজির বক্তব্য হলো, ``দলের সবচেয়ে লম্বা মানুষ'' একটি
\emph{ইমপ্রেডিকেটিভ} সংজ্ঞা; তবু এটি নিঃসন্দেহে গ্রহণযোগ্য।

উত্তরে বলা যেতে পারে, এই ক্ষেত্রে মানুষটিকে
\emph{প্রেডিকেটিভ} উপায়েও চিহ্নিত করা যায়। যেমন, হয়তো
তাঁর দিকে আঙুল দেখালেই চলে। তাহলে ইমপ্রেডিকেটিভ সংজ্ঞার
বিরোধিতা স্পষ্টতই সেইসব সত্তার মধ্যে সীমিত রাখতে হবে, যাদের
\emph{কেবল} ইমপ্রেডিকেটিভ উপায়েই চিহ্নিত করা যায়। কিন্তু
তবুও শুনতে হবে কেন এমন ``অপরিহার্য ইমপ্রেডিকেটিভতা'' এত
খারাপ।\footnote{আরও আলোচনার জন্য \citet{Linnebo2010} দেখো।}

অবশ্য সংজ্ঞাকে যদি কোনো বস্তু \emph{তৈরির} নির্দেশনা হিসেবে
পেতে চাই, তবে ইমপ্রেডিকেটিভ সংজ্ঞা গুরুতর সমস্যার সৃষ্টি করে।
এমন সংজ্ঞা পেলে আমরা সত্যিই দুষ্টচক্রে পড়ব: ইমপ্রেডিকেটিভভাবে
নির্দিষ্ট বস্তুটি তৈরি করতে \emph{আগে} সব বস্তু তৈরি করতে হবে
(তার মধ্যে ওই নির্দিষ্ট বস্তুটিও আছে), কারণ সেটির সংজ্ঞায়
সব বস্তুর উল্লেখ আছে। ফলে বস্তুটি তৈরি করার আগেই সেটি তৈরি
করা লাগবে। কিন্তু এটিও ``অপরিহার্যভাবে ইমপ্রেডিকেটিভ''
সংজ্ঞায়িত সেটের বিরুদ্ধে গুরুতর আপত্তি, কেবল যদি সেটকে
আমাদের \emph{তৈরি করা} বস্তু মনে করি। আর আমরা (সম্ভবত)
তা মনে করি না।

তাই ভালো-মন্দ যা-ই হোক, পরে প্রচলিত হওয়া পদ্ধতিতে
প্রেডিকেটিভতা বা ইমপ্রেডিকেটিভতা নিয়ে কঠোর অবস্থান নেওয়া
হয়নি। তার বদলে এসেছে এখনকার সঞ্চয়ী-পর্যায়ক্রমিক পদ্ধতি।
শেষ পর্যন্ত এর সাহায্যে সেটগুলিকে বিভিন্ন ``পর্যায়ে''
বিন্যস্ত করা যায়—প্রেডিকেটিভ পদ্ধতিতে সত্তাগুলিকে যেমন
সেট$_0$, সেট$_1$, সেট$_2$, \ldots{}-এ বিন্যস্ত করা হয়,
\emph{কিছুটা} সেই রকম—কিন্তু তাদের মধ্যে প্রকারগত কোনো
পার্থক্য ধরে নেওয়া হয় না।

\end{document}
